\section{Reset: an operation that is not a gate}
\hypertarget{reset-semantics}{}

The word \emph{reset} appears in two unrelated places: the compiler's internal
\code{RESET} operation, and three buttons in the Circuit builder. Neither is an
ordinary quantum gate.

\subsection{Reversible gates versus RESET}

Every gate in the palette except the meter is \emph{reversible}: its effect
can always be undone. X, Y, Z, H, CNOT, CCNOT, CZ, CY, SWAP, NOT, AND, NAND,
OR, and XOR undo themselves when applied a second time. S, T, and PHASE are
undone by the same phase with the opposite sign. Mathematically each is a
\emph{unitary} matrix $U$, and $U^\dagger$ undoes it.

\code{RESET} is different. It forces one wire to $\ket0$ whatever it held
before:
\[
\ket0\mapsto\ket0,\qquad \ket1\mapsto\ket0,\qquad
\alpha\ket0+\beta\ket1\mapsto\ket0 .
\]
Two different inputs, $\ket0$ and $\ket1$, give the same output, so nothing
can tell afterwards which one was there. The information has been erased, and
no matrix $U$ can undo that. RESET is therefore not unitary, just like
MEASURE.

\begin{center}
\begin{tabularx}{\textwidth}{>{\bfseries}l X X X}
\toprule
& Reversible gates (X, H, CNOT, AND, \ldots) & MEASURE & RESET\\
\midrule
Unitary? & Yes & No & No\\
Can be undone? & Yes, by the inverse gate & No & No\\
Result depends on chance? & No & Yes & No\\
Written in source? & Yes & Yes & No --- write \code{SET wire 0p}\\
Shown on canvas? & Yes & Yes, as the meter & No (hidden)\\
\bottomrule
\end{tabularx}
\end{center}

\subsection{Where RESET comes from}

You never type RESET. The compiler inserts it when a wire must start at zero:
\begin{itemize}
  \item \code{SET Out 0p} on a workspace wire. (\code{SET Out 1p} and
  \code{SET Out sp} are instead prepared by applying X or H to a wire that
  starts at 0.)
  \item \code{INCREASECYCLE}, which zeroes pending workspace before the next
  stage.
  \item \code{RUNCHILD} and \code{CALL}, which zero each parent output wire
  once before expanding the child.
\end{itemize}
RESET steps are hidden from the canvas and filtered from the
\textbf{Execution log} so that the visible circuit shows only the gates you
wrote. They still run during simulation.

\begin{cautionbox}[Reset a wire only when it is not entangled]
In this simulator, RESET keeps the part of the state in which the wire is
already 0 and rescales it; only if that part is empty does it move the
$\ket1$ amplitudes to $\ket0$. On a fresh wire this is exactly ``prepare
$\ket0$''. On a wire entangled with others it also changes the other wires.
For example, resetting one half of a Bell pair forces the other half to 0 as
well. In multi-stage circuits, \emph{uncompute} temporary wires back to 0 with
reversible gates instead of relying on a reset.
\end{cautionbox}

\subsection{Reset buttons compared}
\hypertarget{reset-buttons}{}

None of these buttons inserts a RESET operation into your circuit. They only
change what the page is showing and simulating.

\begin{center}
\begin{tabularx}{\textwidth}{>{\raggedright\arraybackslash}p{0.2\textwidth}>{\centering\arraybackslash}p{0.13\textwidth}>{\centering\arraybackslash}p{0.13\textwidth}>{\centering\arraybackslash}p{0.13\textwidth}X}
\toprule
Kept or cleared? & \textbf{Reset state} & \textbf{Clear circuit} & \textbf{Reset site} & Notes\\
\midrule
Gates on the canvas & kept & cleared & cleared & \\
Number of wires & kept & kept & back to 3 & \\
Start states (\emph{name} start) & kept & kept & all \code{0p} & \\
Protocol editor text & kept & rewritten & default & Clear circuit writes an
empty canvas protocol into the editor.\\
Compiled PARAMS / RETURNVALS & kept & cleared & cleared & \\
Quantum state and measurements & rewound & rewound & rewound & Back to the
start states; log restarts.\\
Workbench selections & kept & kept & defaults & Gate H, target q0, controls
q1 and q2, phase 90\textdegree.\\
Custom gates and catalog & kept & kept & kept & Nothing removes these except
\textbf{Remove} or closing the browser session.\\
\bottomrule
\end{tabularx}
\end{center}

\begin{beginnerbox}[Which one do I want?]
\begin{description}
  \item[Reset state] ``Run the same experiment again.'' Use it between
  repeated runs of a circuit that measures, to collect several samples.
  \item[Clear circuit] ``Keep my wires, start a new circuit.''
  \item[Reset site] ``Put everything back the way it was when I first opened
  the page.'' The \textbf{Reset site} entry in the navigation menu and
  \textbf{Reset site completely} on the \textbf{More} page do the same thing.
\end{description}
\end{beginnerbox}
